% !TEX root = ../Report.tex
% ======================================================================
\chapter{Conclusions and Future Work}
% ======================================================================

This chapter summarises achievements, critiques the system against its goals,
discusses limitations, and proposes future work.

\section{Conclusions}
GiTrip demonstrates that Git-like version control semantics can be meaningfully
applied to collaborative trip planning. The completed system meets the Must-have
objectives from the Project Specification:
\begin{itemize}
  \item repository-based trip storage with commits and branches,
  \item constraint-aware automatic scheduling integrating routing data,
  \item interactive timeline and map editing with commit saving,
  \item three-way merge with a dedicated conflict-resolution UI.
\end{itemize}

Evaluation results provide objective evidence that users can successfully
generate feasible itineraries and resolve merge conflicts within target times.
Survey feedback confirms perceived novelty and practical value, particularly for
group trips and scenarios requiring careful coordination.

\section{Critical Reflection and Lessons Learned}
Several lessons emerged during development and evaluation:
\begin{itemize}
  \item \textbf{Git concepts are powerful but cognitively expensive.} During
        Evaluation Day, even participants who understood the UI hesitated when
        encountering terms like branch and merge. The Easy mode toggle proved
        essential, but discoverability and mental model support remain ongoing
        work.
  \item \textbf{Structured merge is domain-specific.} Early iterations used
        \texttt{JSON.stringify}-based comparison, which caused subtle bugs when
        key ordering differed between branches. Switching to deep structural
        equality, stable IDs, and explicit deletion semantics resolved these
        issues but required substantial reworking of the merge engine.
  \item \textbf{Constraint scheduling requires explainability.} Feedback from
        testing showed that users trusted the auto-scheduler more when
        adjustments (e.g., opening-hours nudges) were visibly annotated with
        reasons.
  \item \textbf{External dependency management matters.} During development,
        OpenRouteService rate limits and occasional Google Directions failures
        highlighted the need for Haversine-based fallback strategies to keep the
        planner usable under degraded conditions.
\end{itemize}

\section{Limitations}
Key limitations include:
\begin{itemize}
  \item \textbf{Accessibility of interactive components:} maps and drag-based
        timelines remain challenging for keyboard-only and screen-reader users.
  \item \textbf{Merge granularity trade-offs:} while field-level conflicts are
        supported for selected fields, a fully general per-field resolver could
        overwhelm users.
  \item \textbf{No real-time collaboration:} GiTrip supports asynchronous
        collaboration via branches and merges, but not simultaneous live editing
        (CRDT/OT).
  \item \textbf{Heuristic scheduling:} the planner prioritises feasibility and
        responsiveness over global optimality.
\end{itemize}

\section{Future Work}
Future work opportunities are substantial:
\begin{itemize}
  \item \textbf{Improved onboarding and guidance:} in-product tutorials for
        branching/merging, contextual tips, and clearer ``Save'' semantics to
        reduce confusion.
  \item \textbf{Accessibility enhancements:} keyboard-operable alternatives for
        timeline edits, text-based route editing, and stronger ARIA support.
  \item \textbf{Real-time collaboration:} explore CRDT-based shared plan editing,
        while preserving a mergeable history model.
  \item \textbf{Richer merge resolution:} expand conflict detection to additional
        fields, include conflict previews, and support partial merges.
  \item \textbf{Routing extensions:} support ferries/planes as explicit legs,
        integrate GTFS-based transit where applicable, and improve offline routing
        approximations.
  \item \textbf{Optional ML/AI extensions:} while excluded from the MVP, future
        versions could use machine learning to:
        \begin{itemize}
          \item learn user preferences (pace, activity types) and suggest better
                defaults,
          \item predict realistic dwell times and travel-time uncertainty,
          \item recommend destinations based on context (with careful privacy
                design).
        \end{itemize}
        These ideas should be implemented cautiously to maintain explainability
        and avoid undermining deterministic feasibility guarantees.
\end{itemize}

% ======================================================================
\chapter*{Author's Assessment of the Project}
\addcontentsline{toc}{chapter}{Author's Assessment of the Project}
% ======================================================================

\textbf{What is the (technical) contribution of this project?} GiTrip contributes
a practical system that applies Git-style branching, committing, and three-way
merging to structured itinerary data, integrated with constraint-aware scheduling
and interactive timeline/map editing.

\textbf{Why is this contribution relevant to Computer Science?} The project
connects software configuration management concepts (history, merging, conflict
resolution) with human-centred planning, demonstrating how domain semantics
change the design of storage, algorithms, and user interfaces.

\textbf{How can others make use of the work?} Developers can reuse the
snapshot-based repository model, merge engine patterns for structured JSON, and
the planner's constraint-handling approach in other collaborative planning
domains.

\textbf{Why should this project be considered an achievement?} It delivers an
end-to-end deployed application combining routing integration, a scheduler, and a
merge UI; evaluation shows users can complete core tasks successfully with limited
guidance.

\textbf{What are the limitations?} The main limitations are accessibility for
interactive controls, dependence on external routing providers, and the absence
of real-time collaboration and AI-based recommendation within the MVP scope.
